\clearpage
\fancyhead[L]{\small TABLES DE LA LUNE.}
\SourceLeaf{645}
\begin{center}\small Göttingische gelehrte Anzeigen. No. 23. pp. 225..227. 7 February 1825.\\[3pt]\rule{.72\linewidth}{.3pt}\end{center}

\emph{Tables de la lune, formées par la seule théorie de l'attraction, et suivant la division de la circonférence en \textup{400} degrés. Par le Baron \textup{\textsc{de Damoiseau}}, Lieutenant Colonel etc.} Paris, 1824. From Bachelier, 90 pages in 4to.

For all the perfection the lunar tables have attained through the work of \textsc{Tobias Mayer}, \textsc{Bürg} and \textsc{Burckhardt}, their empirical origin remained in some measure a reproach to theoretical astronomy. The accuracy with which the inequalities of the Moon's motion could hitherto be derived from the principle of universal gravitation still fell considerably short of that attained by empirical means. A wreath of honor was still to be won here by tackling very intricate analytical calculations skillfully and carrying them through carefully and tirelessly. The Paris Academy encouraged the effort by setting these calculations as a prize problem, thereby prompting the successful work of \textsc{Damoiseau}, \textsc{Carlini} and \textsc{Plana}. We hope soon to receive the theoretical investigations of these distinguished astronomers. Here we have, to begin with, the tables based on \textsc{Damoiseau}'s derivations. They already attest to the success of those derivations.

The arrangement of these tables differs from that used hitherto in that, for the Moon's longitude in the ecliptic and for the horizontal parallax, all the arguments vary in proportion to time, whereas the latitude is based on the longitude already corrected for the inequalities. This arrangement does considerably reduce convergence and require a greater number of equations (for longitude they are arranged in 46 tables, with contributions from two tables for the perturbations due to Venus and Jupiter). Yet the completely uniform procedure now followed in all the calculations probably offsets this to a considerable extent, although grouping many small equations into double-entry tables still makes them somewhat awkward to use. If double-entry tables had been avoided altogether, the total number of tables for longitude would have risen to 67. Adopting the arrangements chosen by \textsc{Carlini} in his solar tables, where the daily change in each uniformly increasing argument is taken as its unit, would probably make the tables much easier to use.

To make it possible to assess the accuracy attained by \textsc{Damoiseau}'s lunar theory and compare it with that of the earlier tables, a comparison of 50 observations made at Greenwich Observatory over a year and a half in 1802 and 1803 with the tables of \textsc{Bürg}, \textsc{Burckhardt} and \textsc{Damoiseau} is appended at the end of the work. A glance already shows that the last set of tables represents the observations no worse than the earlier ones. To arrive at a more definite judgment, however, the reviewer took the trouble to derive the so-called mean error from the comparison. It worked out as follows.
\begin{center}
\begin{tabular}{l|c|c}
for the tables of & in longitude & in latitude \\
\hline
Bürg & $5''80$ & $9''18$ \\
Burckhardt & $5.37$ & $7.82$ \\
Damoiseau & $5.23$ & $7.12$ \\
\end{tabular}
\end{center}
\noindent So far, then, as this still relatively small number of comparisons permits us to judge, theory has now indeed overtaken empiricism.
Moreover, a not altogether \SourceLeaf{646} inconsiderable part of the discrepancies may well be attributable to the observations themselves.
But now that the way has been opened this far, the hope seems warranted that by proceeding tirelessly along the same path, the tables will be improved ever further and the Moon's motion will be represented with almost the same precision that the refined art of observation can attain.
